\subsection{Arithmetic realization and propagation of ordinary support}
\label{subsec:selmer-arithmetic-propagation}

The normalized curves, auxiliary primes and cohomology systems can now
be constructed with the original arithmetic data.  This supplies the
source's full \texttt{PfromAut} implication, including both
\texttt{P1core} and \texttt{P3core}\cite{SelmerArPropSource}.

\begin{theorem}\label{thm:selmer-arithmetic-propagation}
Retain the original odd prime $\ell$, CM field $L/F$, Hida setting,
base ramification set $S$, coefficient DVR $O$, rank $n>0$, residual
datum, universal ordinary determinant pair $(B,u)$ and ordinary Hecke
support ideal $J$.  Assume the original local class field theory input,
nontrivial residual cyclotomic character, and actual \texttt{HasAut}
module supplier, and take the original base level with $V_0.S'=S$.
Every coefficient curve below maps $u$ to the full determinant and
ordered characters of its stated representation, exactly as in
the original \texttt{IsPt} condition.
Use classical finite local/global duality and Euler
characteristics.  Put $d=[F:\mathbb Q]$.

For every original NiceI coefficient curve
$f:B\to k[[t]]$, if $J\subseteq\ker f$, then
$J\subseteq\mathfrak q$ for every prime $\mathfrak q\subseteq\ker f$.
For every finite chain of original NiceE coefficient curves, suppose
successive kernels contain a common thick prime and the first kernel
contains a stable prime containing $J$.  Then every stable or thick
prime below the last kernel contains $J$.

Here the original distinct ordered ordinary characters, noncriticality,
simple-spectrum density on the full cyclotomic kernel $G'$, and
distinguished-decomposition-group rigidity remain hypotheses of
``nice.''  Stable means absolutely irreducible determinant; thick means
nonstable with nonsemisimple original ordinary equation algebra after
the prescribed algebraic scalar extension.  In particular the theorem
does not replace that closed equation ideal by a reduced or trace ideal.
\end{theorem}

\begin{proof}
\emph{Normalize the actual integral representation.}
Write $A=k[[t]]$ and $K=k((t))$.  The source polarization is exactly
\[
 \rho(g)^tP\rho(cgc^{-1})=\epsilon(g)^{1-n}P.               \tag{RP1}
\]
For NiceI use one block; for NiceE use the two prescribed nonisomorphic
absolutely irreducible $G'$-blocks.  Their given polarizations and
Schur's lemma make the cross blocks of $P$ zero, so its restrictions
are nondegenerate and symmetric.  The commutant consists of one scalar
on each block, so the double-centralizer theorem gives the full rational
matrix algebra, or product of the two block algebras
\cite[IV.1.14--1.16]{SelmerArPropCFT}.

The original integral representation is bounded.  After its fixed
rational block change of basis it remains bounded, so the span of the
translates of a lattice in each block is a stable finite free lattice.
Take their direct sum $M$.  Equation (RP1), its unit multiplier and
the involution $g\mapsto cgc^{-1}$ make the polarization dual $M^\#$
stable too.  Since $2$ is invertible, integral Gram--Schmidt in each
block gives
\[
 P=\operatorname{diag}(t^{a_i}u_i),\qquad
 u_i\in A^\times,\quad a_i\in\mathbb Z.
\]
Indeed choose a diagonal entry of minimum valuation; if there is none,
the sum of the basis vectors for an off-diagonal minimum supplies one.
All Gram--Schmidt subtraction coefficients are integral.  The dual
basis is $t^{-a_i}u_i^{-1}e_i$, so both stabilities give
\[
 v_t(\rho(g)_{ij})\ge\max(0,a_j-a_i).                      \tag{RP2}
\]

Take the quadratic finite-field extension $k_2/k$ and $A'=k_2[[s]]$
with $t=s^2$.  Every nonzero element of $k$ becomes a square in $k_2$;
Hensel's lemma therefore supplies square roots of all $u_i(s^2)$ in
this same ring.  With
\[
 T=\operatorname{diag}(s^{-a_i}u_i(s^2)^{-1/2}),\qquad
 \rho'=T^{-1}\rho T,
\]
we have $T^tPT=1$ and, by (RP2),
\[
 v_s(\rho'(g)_{ij})\ge
 2\max(0,a_j-a_i)+a_i-a_j=|a_i-a_j|\ge0.                  \tag{RP3}
\]
The same holds for $g^{-1}$.  Thus $\rho'$ is integral and has form
matrix $1$, with the original blocks.  The coefficient extension is
finite, injective and local, of rank four.  This is not an index bound
for the two lattices.  Continuity survives fixed conjugation.  Conjugate
$\sum X_g\rho(g)$ over the fraction field and descend its determinant
identity coefficient by coefficient: this preserves the full polynomial
law, including when $\ell\le n$.

The rational ordinary flag respects the blocks.  A block idempotent's
map from a flag step to its quotient is zero because their character
composition factors are disjoint.  Intersect the flag with the normalized
lattice and choose compatible bases of its saturated steps.  Each new
rank-one vector lies in its prescribed block.  These are the exact
integral local bases, zero-entry conditions and ordered diagonal
characters required by NP.  The characters are continuous by rational
triangularization.  They kill the closed commutator and the
$\ell$-primary torsion in the topological abelianization, since $A'^\times$
has no nontrivial $\ell$-power torsion.  Thus the source condition
$\psi_\mu=1$ holds.

Clear denominators in the full rational block span to obtain precisely
the required power-saturated integral span.  Original inertia vanishing,
distinctness and noncriticality survive the injective extension, as does
density of simple spectra.  The equations defining the distinguished
$\epsilon$-twisted intertwiners have the same kernel after extension:
finitely many suffice on the finite matrix space.  The irreducible
case retains $a=n$.  This constructs every field of \texttt{Norm}.

\emph{Construct the arithmetic auxiliary tower.}
For any chosen finite detecting tuple $\sigma_1,\ldots,\sigma_q\in G'$,
use the original open subgroup
\[
 H=\ker\eta\cap\bigcap_i\ker\bar\chi_i\cap\ker\bar\epsilon
      \quad\text{in }\Gamma_0=G_{F,S}.                    \tag{RP4}
\]
Class field theory gives finite $\operatorname{Hom}(U,\mathbb F_\ell)$
for every open subgroup $U$ of $\Gamma_0$ and of each original local
group.  Iterated $\ell$-Frattini kernels give open normal $W_N\subseteq H$
cofinal among neighborhoods of the kernel $\operatorname{PL}(H)$ of
the maximal pro-$\ell$ quotient.  This follows first in every finite
$\ell$-quotient, then by compactness.

Lift $\sigma_j$ to $g_j\in G_L$.  Since $S$ contains all $\ell$ places,
the original cyclotomic compatibility gives
$\epsilon_L(g_j)=\epsilon_{\Gamma_0}(\sigma_j)=1$ in
$\mathbb Z_\ell^\times$, not merely modulo $\ell$.
Impose the finite congruences to
$g_j^{-1}$ modulo $W_N$ and the cyclotomic character modulo $\ell^N$.
They are compatible because $\sigma_j\in G'$.  Apply Chebotarev in a
finite Galois extension over $F$ cutting out these conditions
\cite[V.3.23]{SelmerArPropCFT}.  Select the place over the prime to
obtain the specified Frobenius element, and choose the primes
successively outside $S$ and the preceding choices.  They split in
$L/F$ because the chosen Frobenius restricts trivially to $L$.

For arithmetic Frobenius $\alpha_j(N)$ put
$\phi_j(N)=\alpha_j(N)^{-1}$, the geometric Frobenius in the source.
With the selected prime $w_j(N)$ of $L$ and its prime $v_j(N)$ of $F$,
\[
 \begin{split}
 Q_j(N)&=Nw_j(N)=Nv_j(N)\equiv1\pmod{\ell^N},\\
 h_N(\phi_j(N))\sigma_j^{-1}&\in W_N.
 \end{split}                                               \tag{RP5}
\]
Take $S_N=S\cup\{v_j(N)\}$, the original level with this set,
$\Gamma_N=G_{F,S_N}$, its restriction $h_N$, and the actual tame
pro-$\ell$ inertia generators $\tau_j(N)$.  The exact relation is
\[
 \phi_j(N)^{-1}\tau_j(N)\phi_j(N)=\tau_j(N)^{Q_j(N)}
                         \pmod{\operatorname{PL}(H_N)},
 \qquad H_N=h_N^{-1}H.                                    \tag{RP6}
\]
The kernel of $h_N$ is normally topologically generated by new inertia.
In the pro-$\ell$ quotient the wild and prime-to-$\ell$ parts vanish,
leaving the chosen generators.  Evaluate a character of $H_N$ on their
conjugates by representatives of $\Gamma_N/H_N$.  Characters with the
same values differ by one descending to $H$.  Consequently
\[
 \#\operatorname{Hom}(H_N,\mathbb F_\ell)
 \le \#\operatorname{Hom}(H,\mathbb F_\ell)
                      \ell^{q[\Gamma_0:H]}.                \tag{RP7}
\]
This is the original cardinal bound, with its growing arithmetic index
retained.  The actual restriction maps, places, local embeddings,
inertia generators and (RP4)--(RP7) supply all TWGrp/TWPl fields, also
for $q=0$.  No quantitative prime-size estimate is asserted.

\emph{Realize the two arithmetic cohomology systems.}
Their local indices are precisely the disjoint union of the $\ell$
places, the original old places, $v_0$ and the new primes; this is
bijective with $S_N$.  Every place is split in $L/F$, so the original
decomposition groups over $L$ are the corresponding groups over $F$,
with the same local maps.  The $v_0$ image in $\Gamma_0$ is the original
$D_0$, and the new local groups contain both $\phi_j$ and $\tau_j$.
The local degrees $d_v=[F_v:\mathbb Q_\ell]=[L_w:\mathbb Q_\ell]$
sum to $d$.  Each real embedding gives its actual involution $c_v$,
with $c_v^2=1$ and $\eta(c_v)=\epsilon(c_v)=-1$.

Finite local duality and the Euler formula give
\[
 \#Z^1(G_v,M)=|M|^{1+d_v}|(M^\vee(1))^{G_v}|,
\]
where $d_v=0$ away from $\ell$.  A perfect source pairing into
$\mathbb F_\ell$ forces $\ell M=0$ and identifies its second module
with $M^\vee(1)$, so this covers every original local-count input.
The global Selmer ratio gives exactly the established first-system
Greenberg--Wiles inequality with all these places and real factors
\cite[I.2.1, I.2.3, I.2.8, I.4.10, I.4.15, I.5.1]{SelmerArPropADT}.

At an old place retain
$J_u=I_u\cap\bigcap_i\ker\bar\chi_i$, whose quotient $I_u/J_u$
is finite of order prime to $\ell$.  For $M$ trivial on $J_u$, averaging
over this quotient and the procyclic unramified group give
\[
 H^1(G_u/J_u,M)=M^{I_u/J_u}/(\mathrm{Frob}-1)M^{I_u/J_u},
 \qquad \#H^1(G_u/J_u,M)=|M^{G_u}|.
\]
Hence there are exactly $|M|$ cocycles vanishing on $J_u$.
This quotient has $\ell$-cohomological dimension at most one.
For any prescribed module surjection, first lift an unramified class
and then its correcting boundary primitive.  This lifts the specified
cocycle itself and proves the entire original \texttt{UrAx} input.

Local and global $H^1,H^2$ are finite.  Their $H^3$ vanishes for these
finite $\ell$-modules: locally by cohomological dimension two, globally
by the real-place description and odd $\ell$.  Thus $H^2$ is right exact
for the original coefficient surjections.  The exact global Euler
formula is
\[
 |H^0(M)|\,|H^2(M)|\,|M|^d
       =|H^1(M)|\prod_{v\text{ real}}|M^{c_v}|.             \tag{RP8}
\]
For the source obstruction cokernel the contract explicitly requires
$\ell M=0$, as holds for the actual characteristic-$\ell$ truncations.
Local right exactness maps the
full $S_N$ sum of local $H^2(M)$ onto its prescribed local-quotient
$H^2$ groups, compatibly with global restriction.  Its cokernel is a
quotient of the full localization cokernel.  The Poitou--Tate tail
therefore gives
\[
             |\operatorname{Ck2}|
                 \le |\operatorname{Hom}_G(M,\mathbb F_\ell(1))|.
\]
This proves every ArSys and CohAx field with the original quotients
and cyclotomic action.  The established B1/B2 theorems now apply to
these actual systems; B1 has exactly the exponent $nd+nq$.
Local reciprocity supplies the retained local class field input from
the finite unit torsion, the free $\mathbb Z_\ell^{[F_v:\mathbb Q_\ell]}$
unit factor and the one valuation factor\cite[I.1.1]{SelmerArPropCFT}.

\emph{Apply the actual marked propagation.}
Normalize every curve as above.  Established finite simultaneous
detection and coefficient extension provide one detecting tuple,
the original odd families, distinct split auxiliary roots and equal
residue-field cardinalities.  Each coefficient map is injective;
all original kernels in $B$ remain unchanged.  Construct the tower,
the two arithmetic systems and their B1/B2 bounds as just proved.
On $H$ every residual matrix is unipotent, so its finite image is an
$\ell$-group by Cauchy's theorem.  Higher congruence kernels are
$\ell$-groups too; hence each point and its finite cocycle actions
factor through the required pro-$\ell$ quotient.

HasAut supplies the actual TWAut modules.  The universal ring is
generated topologically by its full determinant coefficients and
ordered characters: construct their universal rings and quotient by
the original Pair equations; uniqueness identifies this quotient
with $B$.  The image of the old $J_u$ is exactly $r.Ip_u$.
Every hypothesis of the proved Raw theorem is therefore satisfied.

Stable is equivalent to the explicit nonzero $n^2$ by $n^2$ trace-Gram
discriminant.  Indeed the trace pairing on $M_n$ is nondegenerate in
every characteristic; a reducible semisimple representation has image
algebra of dimension less than $n^2$.  Thick implies NonSS directly
for the same original equation algebra.

At a NiceI point the rational group span already gives a nonzero
trace-Gram determinant outside its kernel and every smaller prime.
Apply Raw with a zero-length chain and that kernel as the initial
stable prime containing $J$; this proves P1core.  At a NiceE chain,
the original thick links and starting stable prime are exactly the
Raw hypotheses.  Its conclusion is P3core, whose zero-length case is
P2core.  All ranks, ordered characters and exponent losses are retained.
\end{proof}

The next obligation is the global component argument that reaches the
required characteristic-prime support from the actual seed.  The
theorem does not infer that starting support from a characteristic-zero
point, or supply the final uniform integral Selmer budget.
